\documentclass[9pt]{extarticle}
\usepackage[utf8]{inputenc}
\usepackage[T1]{fontenc}
\usepackage{mathpazo}
\usepackage[a4paper,left=1.1cm,right=1.1cm,top=1.2cm,bottom=1.3cm]{geometry}
\usepackage{amsmath,amssymb}
\usepackage[table]{xcolor}
\usepackage{multicol,tabularx,array}
\usepackage[most]{tcolorbox}
\usepackage{titlesec,enumitem,fancyhdr}
\renewcommand{\sfdefault}{phv}
\definecolor{rosso}{HTML}{D42A2F}
\definecolor{verde}{HTML}{DCEFD7}
\definecolor{verdeb}{HTML}{6AA84F}
\definecolor{giallo}{HTML}{FFF6D5}
\definecolor{gialloc}{HTML}{F5D76E}
\definecolor{azz}{HTML}{E6EEF8}
\definecolor{azzc}{HTML}{3D6FB6}
\definecolor{rosa}{HTML}{FBE9EA}
% sezioni
\newcommand{\secbox}[1]{\colorbox{rosso}{\parbox{\dimexpr\linewidth-2\fboxsep}{\strut\color{white}\thesection\quad #1}}}
\titleformat{\section}{\sffamily\bfseries}{}{0pt}{\secbox}
\titlespacing*{\section}{0pt}{6pt}{3pt}
\titleformat{\subsection}{\sffamily\color{rosso}}{\thesubsection}{0.5em}{}
\titlespacing*{\subsection}{0pt}{4pt}{2pt}
\renewcommand{\thesection}{\arabic{section}}
\setlength{\parindent}{0pt}\setlength{\parskip}{1.5pt}
\setlength{\columnsep}{12pt}\setlength{\columnseprule}{0.4pt}
\setlist{nosep,leftmargin=11pt}
\setlength{\abovedisplayskip}{2pt}\setlength{\belowdisplayskip}{2pt}
\setlength{\abovedisplayshortskip}{1pt}\setlength{\belowdisplayshortskip}{1pt}
\renewcommand{\arraystretch}{1.12}
\setlength{\tabcolsep}{3pt}
% riquadri
\newtcolorbox{bx}[3]{enhanced,breakable,colback=#1,colframe=#1,boxrule=0pt,arc=0pt,left=4pt,right=4pt,top=2pt,bottom=2pt,before skip=4pt,after skip=4pt,
 title={\textsc{#3}},coltitle=#2,colbacktitle=#1,fonttitle=\sffamily\footnotesize,attach title to upper={\par}}
\newenvironment{esame}{\begin{bx}{giallo}{rosso!80!black}{nell'esame}}{\end{bx}}
\newenvironment{svolto}[1]{\begin{bx}{azz}{azzc}{esempio svolto -- #1}}{\end{bx}}
\newenvironment{errori}{\begin{bx}{rosa}{rosso}{errori da evitare}}{\end{bx}}
\newenvironment{oro}[1]{\begin{bx}{giallo}{rosso!80!black}{#1}}{\end{bx}}
% tabelle
\newcolumntype{P}[1]{>{\raggedright\arraybackslash}p{#1}}
\newcolumntype{L}{>{\raggedright\arraybackslash}X}
\newcommand{\hd}{\rowcolor{verde}}
\newenvironment{tab}[1]{\par\noindent\rowcolors{2}{white}{giallo!70}\tabularx{\linewidth}{|#1|}\hline}{\hline\endtabularx\par\vspace{2pt}}
\pagestyle{fancy}\fancyhf{}\renewcommand{\headrulewidth}{0pt}
\fancyfoot[C]{\scriptsize Formulario Fisica -- Enrico Baldasso -- pag. \thepage}
\begin{document}

{\sffamily\LARGE\color{rosso} FORMULARIO --- FISICA}\hfill{\sffamily\bfseries Enrico Baldasso 4685}\\[-1pt]
{\scriptsize\itshape Esame unico Fisica I + Fisica II (4 esercizi F1a, F1b, F2a, F2b + 4 domande D1a, D1b, D2a, D2b). Formule, scalette e casistiche tratte dagli esercizi della dispensa e dal compito d'esempio. Unità SI.}
\vspace{2pt}

\begin{multicols*}{2}
\setcounter{section}{-1}
\section{Come muoversi all'esame: mappa e strategia}
\subsection{Riconosci l'esercizio dalle parole del testo}
\begin{tab}{P{3.3cm}|P{0.5cm}|L}
\hd Se il testo dice\dots & \S & Metodo in una riga\\\hline
``disco/asta che ruota'', ``punto fissato al bordo'', ``piccole oscillazioni'' & \S5 & energia con $\frac12I\omega^2$; pendolo fisico $T=2\pi\sqrt{I/mgd}$\\
``rotola senza strisciare'' & \S5 & $v=\omega R$, \ $K=\frac12(M+I/R^2)v^2$\\
``piano inclinato'', ``carrucola'', ``tensione'' & \S3 & $\sum\vec F=m\vec a$ su ogni corpo\\
``velocità in B'', ``quota massima'', ``molla compressa'' & \S3 & conservazione dell'energia\\
``urto'', ``restano uniti'', ``blocco mobile'' & \S3 & $\sum\vec p$ costante (+ $E$ se elastico)\\
``scala'', ``asta in equilibrio'' & \S5 & $\sum F=0$, $\sum\tau=0$\\
``galleggia'', ``immerso'', ``tubo'', ``strozzatura'' & \S6 & Stevino, Archimede, Bernoulli\\
``ciclo'', ``rendimento'', ``adiabatica'' & \S8 & tabella trasformazioni, $\eta=1-|Q_c|/Q_a$\\
``sorgenti'', ``reversibile?'' & \S8 & Clausius $\sum Q/T\le0$\\
``cariche nei vertici'', ``lavoro per spostare'' & \S9 & $\vec E$ vettoriale, $V$ scalare\\
``sfere concentriche'', ``densità di carica'' & \S9 & Gauss\\
``condensatore'', ``dielettrico'', ``energia elettrostatica'' & \S10 & serie/parallelo, $U=\frac12CV^2$\\
``generatore'', ``potenza massima'' & \S11 & $i=\mathcal E/(R+r)$, $R=r$\\
``ione'', ``raggio'', ``periodo'' in $B$ & \S12 & $r=mv/qB$, $T=2\pi m/qB$\\
``bobina'', ``$B(t)$'', ``f.e.m.'', ``barra che scorre'' & \S13 & $\mathcal E=-d\Phi/dt$, $i=\mathcal E/R$\\
``frange'', ``fenditure'', ``prisma'', ``lastra'' & \S14 & $d\sin\theta=m\lambda$, Snell\\
\end{tab}

\begin{oro}{regole d'oro valide in ogni esercizio}
\begin{itemize}
\item \textbf{Prima la formula letterale}, poi i numeri, poi le unità: il procedimento vale più del numero.
\item \textbf{Scrivi il principio in una frase} (``si conserva l'energia perché non c'è attrito'').
\item \textbf{Dati in SI}: cm$\to$m, km/h$\to$m/s ($\div3{,}6$), $^\circ$C$\to$K, cal$\to$J, diametro$\to$raggio.
\item Se il testo è ambiguo (raggio/diametro, $H$ o $B$) \textbf{dichiara l'interpretazione}.
\end{itemize}
\end{oro}

\begin{bx}{azz}{azzc}{piano di lavoro suggerito}
\begin{enumerate}[label=\textcolor{rosso}{\bfseries\arabic*.}]
\item 5 min: leggi tutto e riconosci i 4 esercizi con la tabella 0.1.
\item Domande teoriche per prime: scaletta \S15 (immagine $\to$ definizione $\to$ formula $\to$ esempio).
\item Esercizi dal più sicuro; se blocchi oltre 5 min scrivi la scaletta e passa oltre.
\item Ultimi 10 min: unità, segni, kelvin, angoli, ordini di grandezza.
\end{enumerate}
\end{bx}

\section{Strumenti di base}
\begin{tab}{L|L}
\hd Costante & Valore\\\hline
$g$ & $9{,}81$ m/s$^2$\\
$R$ gas, $k_B$ & $8{,}314$ J/(mol\,K), \ $1{,}38\cdot10^{-23}$ J/K\\
$k=\frac1{4\pi\varepsilon_0}$, \ $\varepsilon_0$ & $8{,}99\cdot10^9$, \ $8{,}854\cdot10^{-12}$ F/m\\
$\mu_0$ & $4\pi\cdot10^{-7}$ H/m\\
$e$, $m_e$, $m_p$ & $1{,}6\cdot10^{-19}$ C, $9{,}11\cdot10^{-31}$, $1{,}67\cdot10^{-27}$ kg\\
$c$, $h$ & $3\cdot10^8$ m/s, \ $6{,}63\cdot10^{-34}$ J\,s\\
1 atm, 1 cal, 1 eV & $1{,}013\cdot10^5$ Pa $=760$ mmHg; $4{,}186$ J; $1{,}6\cdot10^{-19}$ J\\
giri/min $\to$ rad/s & $\times2\pi/60$\\
\end{tab}
Vettori: $A_x=A\cos\theta$, $A_y=A\sin\theta$, $|\vec A|=\sqrt{A_x^2+A_y^2}$; $\vec a\cdot\vec b=ab\cos\theta$, $|\vec a\times\vec b|=ab\sin\theta$.\\
Grafici: pendenza $=$ derivata, area $=$ integrale ($F=-dU/dx$, $W=$ area sotto $p(V)$).

\section{Cinematica}
\begin{tab}{P{2.3cm}|L}
\hd Moto & Equazioni\\\hline
Rettilineo unif. & $x=x_0+vt$\\
Unif. accelerato & $x=x_0+v_0t+\frac12at^2$, \ $v=v_0+at$, \ $v^2=v_0^2+2a\Delta x$\\
Lancio verticale & $h_{max}=\frac{v_0^2}{2g}$, \ $t_{sal}=\frac{v_0}{g}$\\
Parabolico & $x=v_0\cos\theta\,t$, \ $y=v_0\sin\theta\,t-\frac12gt^2$\\
& traiettoria $y=x\tan\theta-\frac{gx^2}{2v_0^2\cos^2\theta}$; gittata $\frac{v_0^2\sin2\theta}{g}$\\
Circolare & $v=\omega R$, \ $a_c=\omega^2R=\frac{v^2}{R}$, \ $\omega=\frac{2\pi}{T}=2\pi f$\\
Relativo & $\vec v_{ass}=\vec v_{rel}+\vec v_{trasc}$ (ombrello, pioggia)\\
\end{tab}
\begin{svolto}{9.5}
Lancio a $60^\circ$, bersaglio a $x=24$, $y=13{,}4$ m. Dalla traiettoria: $v_0^2=\frac{gx^2}{2\cos^2\theta(x\tan\theta-y)}$ $\Rightarrow$ $v_0=20{,}0$ m/s; $t=x/v_x=2{,}40$ s, $v_y=v_0\sin\theta-gt=-6{,}2$ $\Rightarrow$ $v=11{,}76$ m/s a $31^\circ36'$.
\end{svolto}

\section{Dinamica, lavoro, energia, urti}
\begin{tab}{P{2.3cm}|L}
\hd Legge & Formula\\\hline
Newton & $\sum\vec F=m\vec a=\frac{d\vec p}{dt}$\\
Forze & $P=mg$, \ $f_s\le\mu_sN$, \ $f_d=\mu_dN$, \ $F_{el}=-kx$\\
Piano inclinato & $N=mg\cos\alpha$, \ $a=g(\sin\alpha-\mu\cos\alpha)$\\
Carrucola ideale & $a=\frac{m_Bg}{m_A+m_B}$, \ $T=m_Aa$\\
Curva & $v_{max}=\sqrt{\mu gr}$; \ pendolo nel punto basso $T=mg+\frac{mv^2}{l}$\\
Lavoro, potenza & $W=\int\vec F\cdot d\vec s$, \ $W_{tot}=\Delta K$, \ $P=Fv$\\
Energia & $K=\frac12mv^2$, \ $U_g=mgh$, \ $U_{el}=\frac12kx^2$\\
Conservazione & $K_i+U_i=K_f+U_f$; \ con attrito $\Delta E=-\mu Nd$\\
Urto anelastico & $v'=\frac{m_1v_1+m_2v_2}{m_1+m_2}$\\
Urto elastico 1D & $v_1'=\frac{(m_1-m_2)v_1+2m_2v_2}{m_1+m_2}$; masse uguali: si scambiano\\
Centro di massa & $\vec F_{est}=M\vec a_{cm}$; \ se $F_{est,x}=0$: $m_1x_1+m_2x_2=$ cost.\\
Gravitazione & $F=G\frac{Mm}{r^2}$, \ orbita $v=\sqrt{GM/r}$\\
\end{tab}
\begin{svolto}{11.4}
Punto $m$ che scivola su blocco $M$ libero: $mv=MV$ e $mgR=\frac12mv^2+\frac12MV^2$ $\Rightarrow$ $v=\sqrt{\frac{2gRM}{M+m}}$, $V=\frac mMv$.
\end{svolto}

\section{Oscillazioni e onde}
\begin{tab}{P{2.3cm}|L}
\hd Sistema & Formula\\\hline
Molla & $\omega=\sqrt{k/m}$, \ $T=2\pi\sqrt{m/k}$, \ $x=A\cos(\omega t+\phi)$\\
Due masse + molla & massa ridotta $\mu=\frac{m_1m_2}{m_1+m_2}$, \ $\omega=\sqrt{k/\mu}$\\
Pendolo semplice & $T=2\pi\sqrt{l/g}$\\
Pendolo fisico & $T=2\pi\sqrt{\frac{I_O}{m_{tot}gd}}$ \ ($d$ perno--CM)\\
Onde & $v=\lambda f$, \ corda $v=\sqrt{T/\mu}$, \ suono $340$ m/s\\
Stazionarie & estremi fissi $f_n=\frac{nv}{2L}$\\
Pozzo / eco & $t=\sqrt{2h/g}+h/v_s$\\
\end{tab}

\section{Corpo rigido e statica}
\begin{tab}{P{2.6cm}|L}
\hd Corpo (asse per CM) & $I$\\\hline
punto a distanza $r$ & $mr^2$\\
anello / cilindro cavo & $MR^2$\\
disco / cilindro pieno & $\frac12MR^2$\\
sfera piena / cava & $\frac25MR^2$ \ / \ $\frac23MR^2$\\
asta (centro / estremo) & $\frac1{12}ML^2$ \ / \ $\frac13ML^2$\\
\end{tab}
Steiner $I=I_{cm}+Md^2$ \quad $\tau=rF\sin\theta=I\alpha$ \quad $K=\frac12I\omega^2$ \quad $L=I\omega$ (cost. se $\tau_{est}=0$)\\
\textbf{Rotolamento:} $v=\omega R$; \ $K=\frac12(M+\frac{I}{R^2})v^2$; \ da quota $h$: $v=\sqrt{\frac{2gh}{1+I/MR^2}}$ (cilindro $\sqrt{4gh/3}$); \ $a=\frac{g\sin\alpha}{1+I/MR^2}$.\\
Striscia e poi rotola: $v=v_0-\mu gt$, $\omega=\omega_0+\frac{\mu MgR}{I}t$, fine per $v=\omega R$.\\
\textbf{Statica:} $\sum\vec F=0$, $\sum\vec\tau=0$ rispetto al polo con più incognite. Scala su parete liscia: $N_s=P_{tot}$, $N_p=f\le\mu N_s$.

\subsection{Scalette F1a -- corpo che ruota attorno a un asse fisso}
\begin{tab}{c|P{2.9cm}|L}
\hd N & Cosa fai & Scrivi\\\hline
1 & Disegno: posizioni iniziale e finale & quanto scende ogni massa ($\Delta h$)\\
2 & Momento d'inerzia totale & $I=I_{corpo}+\sum m_ir_i^2$ (+ Steiner)\\
3 & Conservazione dell'energia & $\sum m_ig\Delta h_i=\frac12I\omega^2$, $v=\omega r$\\
4 & Piccole oscillazioni & $I\ddot\theta=-mgd\,\theta$ $\Rightarrow$ $\omega_0^2=\frac{mgd}{I}$\\
5 & Controllo & caso limite ($M\to0$: pendolo semplice)\\
\end{tab}
\begin{svolto}{F1a compito}
Disco $M,R$ + massa $m$ al bordo, da orizzontale al punto basso: $mgR=\frac12(\frac12M+m)v^2$ $\Rightarrow$ $v=\sqrt{\frac{4mgR}{M+2m}}$; \ $T=2\pi\sqrt{\frac{(M+2m)R}{2mg}}$.
\end{svolto}

\section{Fluidi}
\begin{tab}{P{2.3cm}|L}
\hd Legge & Formula\\\hline
Stevino & $p=p_0+\rho gh$ \ (acqua $10^3$, Hg $13{,}6\cdot10^3$ kg/m$^3$)\\
Pascal (torchio) & $\frac{F_1}{A_1}=\frac{F_2}{A_2}$\\
Archimede & $S=\rho_{fl}V_{imm}g$; \ galleggia: $\frac{V_{imm}}{V}=\frac{\rho_c}{\rho_{fl}}$\\
Peso apparente & $P'=P(1-\rho_{fl}/\rho_c)$\\
Continuità & $Q=A_1v_1=A_2v_2$\\
Bernoulli & $p+\frac12\rho v^2+\rho gz=$ cost.; \ Torricelli $v=\sqrt{2gh}$\\
Gas in tubo/campana & Boyle $p_1V_1=p_2V_2$ + Stevino\\
\end{tab}

\section{Calore e gas perfetti}
\begin{tab}{P{2.3cm}|L}
\hd Legge & Formula\\\hline
Calorimetria & $Q=mc\Delta T$; \ $\sum m_ic_i(T_e-T_i)=0$; \ $c_{acqua}=1$ cal/g$^\circ$C\\
Cinetica$\to$calore & $\frac12mv^2=mc\Delta T$\\
Dilatazione & $\Delta L=\lambda L_0\Delta T$, \ $\Delta V=\beta V_0\Delta T$ ($\beta\simeq3\lambda$)\\
Gas perfetto & $pV=nRT$; \ $n$ cost.: $\frac{p_1V_1}{T_1}=\frac{p_2V_2}{T_2}$; \ $m\propto\frac{pV}{T}$\\
Calori specifici & mono $c_V=\frac32R$, $\gamma=\frac53$; \ bi $c_V=\frac52R$, $c_p=\frac72R$, $\gamma=1{,}4$\\
Scale & $T_K=T_C+273{,}15$, \ $T_F=\frac95T_C+32$\\
\end{tab}

\section{Termodinamica}
\begin{tab}{P{1.4cm}|P{2.1cm}|P{1.3cm}|L}
\hd Trasf. & $W$ & $Q$ & note\\\hline
isocora & $0$ & $nc_V\Delta T$ & $p/T$ cost.\\
isobara & $p\Delta V$ & $nc_p\Delta T$ & $V/T$ cost.\\
isoterma & $nRT\ln\frac{V_f}{V_i}$ & $=W$ & $\Delta U=0$\\
adiabatica & $-\Delta U=\frac{p_iV_i-p_fV_f}{\gamma-1}$ & $0$ & $pV^\gamma$, $TV^{\gamma-1}$ cost.\\
\end{tab}
$\Delta U=Q-W$ (sempre $\Delta U=nc_V\Delta T$) \quad ciclo: $\Delta U=0$, $W=Q_a-|Q_c|$\\
$\eta=\frac{W}{Q_a}=1-\frac{|Q_c|}{Q_a}$ \quad Carnot $\eta=1-\frac{T_f}{T_c}$ \quad Frigo $COP=\frac{Q_f}{W}$\\
Clausius $\sum\frac{Q_i}{T_i}\le0$ ($=0$ reversibile) \quad $\Delta S=nc_V\ln\frac{T_f}{T_i}+nR\ln\frac{V_f}{V_i}$\\
Macchine in cascata: $\eta_{tot}=1-(1-\eta_1)(1-\eta_2)$

\subsection{Scalette F1b -- ciclo con gas perfetto}
\begin{tab}{c|P{2.9cm}|L}
\hd N & Cosa fai & Scrivi\\\hline
1 & Temperature in kelvin & $T=T_C+273{,}15$\\
2 & Per ogni tratto: segno di $Q$ & dove $T$ sale a $p$ o $V$ cost. $\to$ $Q_{ass}$\\
3 & Temperature mancanti & adiabatiche $TV^{\gamma-1}$ cost., isobare $V\propto T$, isocore $p\propto T$\\
4 & Calori & $Q_a=\sum nc\,\Delta T$, \ $|Q_c|$ idem\\
5 & Rendimento e controllo & $\eta=1-\frac{|Q_c|}{Q_a}<\eta_{Carnot}$\\
\end{tab}
\begin{svolto}{F1b compito}
BC adiab., CD isobara, DE adiab., EB isocora, biatomico. $T_E=T_B(T_D/T_C)^\gamma=1051$ K; \ $\eta=1-\frac1\gamma\frac{T_E-T_B}{T_D-T_C}=1-\frac{737{,}6}{1{,}4\cdot1200}=\mathbf{0{,}56}$.
\end{svolto}
\begin{svolto}{24.1}
4 sorgenti: $\sum Q/T=\frac{68}{573}-\frac{14}{353}+\frac{22}{298}-\frac{51}{203}<0$ $\Rightarrow$ irreversibile; $\eta=\frac{68-14+22-51}{68+22}=0{,}28$.
\end{svolto}
\begin{errori}
$^\circ$C invece di K \textbullet\ rendimento diviso per il calore \emph{totale} anziché quello \emph{assorbito} \textbullet\ $c_p$ nell'isocora o $c_V$ nell'isobara \textbullet\ $\gamma$ sbagliato (mono/bi).
\end{errori}

\section{Elettrostatica}
\begin{tab}{P{2.6cm}|L}
\hd Legge / caso & Formula\\\hline
Coulomb & $F=k\frac{q_1q_2}{r^2}$, \ $\vec E=\vec F/q$, \ $E=k\frac{q}{r^2}$\\
Potenziale & $V=k\frac qr$, \ $\Delta V=-\int\vec E\cdot d\vec l$, \ $W_{est}=q\Delta V$, \ $E$ unif. $\Delta V=Ed$\\
Gauss & $\oint\vec E\cdot d\vec A=\frac{Q_{int}}{\varepsilon_0}$\\
sfera/guscio ($r>R$) & $E=k\frac{Q}{r^2}$; \ dentro guscio $0$; sfera piena $E=\frac{kQr}{R^3}$\\
filo / piano & $E=\frac{\lambda}{2\pi\varepsilon_0r}$ \ / \ $E=\frac{\sigma}{2\varepsilon_0}$\\
conduttore & dentro $E=0$; fuori $E=\frac{\sigma}{\varepsilon_0}$\\
asse di un disco & $E=\frac{\sigma}{2\varepsilon_0}\big(1-\frac{z}{\sqrt{z^2+R^2}}\big)$\\
dipolo & $p=qd$, \ $\tau=pE\sin\theta$\\
\end{tab}
\begin{svolto}{28.1}
Due $q$ uguali a 12 cm: nel punto medio $E=0$; sull'asse a 6 cm $E=2\frac{kq}{r^2}\cos45^\circ$ (componenti $\parallel$ AB si annullano); $\Delta V=2kq(\frac1{r_D}-\frac1{r_O})$.
\end{svolto}
\begin{errori}
Sommare i moduli di $\vec E$ invece delle componenti \textbullet\ distanza dal \emph{centro} del quadrato $=l/\sqrt2$.
\end{errori}

\section{Condensatori e dielettrici}
\begin{tab}{P{2.6cm}|L}
\hd Caso & Formula\\\hline
piano & $C=\frac{\varepsilon_0\varepsilon_rA}{d}$, \ $E=\frac Vd$\\
sferico / cilindrico & $C=4\pi\varepsilon_0\varepsilon_r\frac{r_1r_2}{r_2-r_1}$ \ / \ $\frac{2\pi\varepsilon_0\varepsilon_rL}{\ln(r_2/r_1)}$\\
parallelo (stessa $V$) & $C=C_1+C_2$ \ (dielettrico su \emph{parte dell'area})\\
serie (stessa $Q$) & $\frac1C=\frac1{C_1}+\frac1{C_2}$ \ (dielettrico su \emph{parte dello spessore})\\
energia & $U=\frac12CV^2=\frac{Q^2}{2C}$, \ $u=\frac12\varepsilon_0\varepsilon_rE^2$\\
dielettrico & $C=\varepsilon_rC_0$; \ $Q$ fissa: $E=E_0/\varepsilon_r$; \ $V$ fissa: $Q$ aumenta\\
conduttori collegati & $V=\frac{Q_A+Q_B}{C_A+C_B}$, energia persa $U_i-U_f$\\
interfaccia & $\frac{\tan\alpha_1}{\tan\alpha_2}=\frac{\varepsilon_1}{\varepsilon_2}$, $E_t$ e $D_n$ continui\\
\end{tab}
\subsection{Scalette F2a -- condensatore composto}
\begin{tab}{c|P{2.9cm}|L}
\hd N & Cosa fai & Scrivi\\\hline
1 & Disegno in sezione & zone con dielettrico/distanza diversi\\
2 & Stessa $V$ o stessa $Q$? & parallelo (affiancate) / serie (sovrapposte)\\
3 & Aree e distanze & raggio vs diametro! corona $\pi(r_2^2-r_1^2)$\\
4 & $C$ totale, ricava l'incognita & formula letterale prima dei numeri\\
5 & Energia & $U=\frac12CV^2$\\
\end{tab}
\begin{svolto}{F2a compito}
Disco ($\varepsilon_r=10$, $d_1$) + corona (vuoto, $d_1/2$) in parallelo: $C=\frac{\varepsilon_0}{d_1}(\varepsilon_rA_1+2A_2)$ $\Rightarrow$ $d_1=1{,}36$ mm; \ $U=\frac12\cdot25\,\text{nF}\cdot100^2=1{,}25\cdot10^{-4}$ J.
\end{svolto}

\section{Corrente elettrica}
$i=\frac{dq}{dt}$, \ $V=Ri$, \ $R=\rho\frac{L}{A}$, \ $P=Vi=Ri^2=\frac{V^2}{R}$\\
Generatore reale: $i=\frac{\mathcal E}{R+r}$, $P_R=\frac{\mathcal E^2R}{(R+r)^2}$, massima per $R=r$: $\frac{\mathcal E^2}{4r}$\\
Serie $\sum R_i$; parallelo $\frac1R=\sum\frac1{R_i}$; Kirchhoff nodi e maglie\\
RC: $q=Q_0e^{-t/RC}$ \quad RL all'apertura: $i=I_0e^{-t(R+r)/L}$, carica $Q=\frac{LI_0}{R+r}$\\
Riscaldamento: $P\,\Delta t=mc\,\Delta T$

\section{Campo magnetico}
\begin{tab}{P{2.6cm}|L}
\hd Caso & Formula\\\hline
Lorentz & $\vec F=q\vec v\times\vec B$; \ $r=\frac{mv}{qB}=\frac{\sqrt{2mE_k}}{qB}$, \ $T=\frac{2\pi m}{qB}$\\
filo in $B$ & $F=iLB\sin\theta$\\
spira in $B$ & $m=NiA$, \ $\tau=mB\sin\theta$ \ (max se piano $\parallel$ $B$)\\
filo infinito & $B=\frac{\mu_0i}{2\pi r}$\\
centro spira & $B=\frac{\mu_0Ni}{2R}$; \ asse $\frac{\mu_0iR^2}{2(R^2+z^2)^{3/2}}$\\
tratto finito & $B=\frac{\mu_0i}{4\pi d}(\sin\theta_1+\sin\theta_2)$; quadrato $\frac{2\sqrt2\mu_0i}{\pi a}$\\
solenoide / toroide & $B=\mu_0ni$ \ / \ $\frac{\mu_0Ni}{2\pi r}$\\
nastro largo $l$ & $H=\frac{i}{2\pi l}\ln\frac{s+l}{s}$\\
fili paralleli & $\frac FL=\frac{\mu_0i_1i_2}{2\pi d}$\\
materia & $B=\mu_0\mu_rH$; \ $NI=\Phi\mathcal R$, \ $\mathcal R=\frac{l}{\mu_0\mu_rA}$\\
\end{tab}
\begin{errori}
\textbf{A/m (Asp/m) è $H$, non $B$}: $B=\mu_0H$ \textbullet\ $m_p=1{,}66\cdot10^{-27}$ kg (non $10^{-17}$).
\end{errori}

\section{Induzione elettromagnetica}
$\Phi=NBA\cos\theta$ \ ($\theta$ tra $\vec B$ e la \textbf{normale}; se l'angolo è col \emph{piano}: $\sin$)\\
$\mathcal E=-\frac{d\Phi}{dt}$, \ $i=\frac{\mathcal E}{R}$, \ $P=Ri^2$, \ carica $q=\frac{\Delta\Phi}{R}$\\
Barra: $\mathcal E=Blv$ \quad Disco di Barlow: $\mathcal E=\frac12B\omega R^2$ \quad Spira rotante: $NBA\omega\sin\omega t$\\
$B=B_0e^{-t/\tau}$ $\Rightarrow$ $\mathcal E=\frac{NAB_0\cos\theta}{\tau}e^{-t/\tau}$\\
$L=\mu_0n^2Al$, \ $\mathcal E=-L\frac{di}{dt}$, \ $U=\frac12Li^2$, \ $u=\frac{B^2}{2\mu_0}$; \ filo: $\frac Ul=\frac{\mu_0i^2}{16\pi}$\\
Mutua: $\Phi_2=Mi_1$; forza su piccola bobina $F=m\frac{dB}{dx}$
\subsection{Scalette F2b -- bobina in campo variabile}
\begin{tab}{c|P{2.9cm}|L}
\hd N & Cosa fai & Scrivi\\\hline
1 & Area e angolo & $A=\pi r^2$ ($r=d/2$); angolo con la normale\\
2 & Flusso & $\Phi(t)=NB(t)A\cos\theta$\\
3 & f.e.m. & $\mathcal E=-d\Phi/dt$ (deriva $B(t)$)\\
4 & Corrente e incognita & $i=\mathcal E/R$; risolvi per $t$ con $\ln$\\
5 & Potenza e controllo & $P=Ri^2$; verifica che $i_{max}$ ($t=0$) basti\\
\end{tab}
\begin{svolto}{F2b compito}
200 spire, $d=10$ cm, $R=100\,\Omega$, $30^\circ$ col piano: $\mathcal E=\frac{NA\sin30^\circ B_0}{\tau}e^{-t/\tau}$; $i=1$ A $\Rightarrow$ $t=\tau\ln\frac{\mathcal E_0}{100}\simeq14{,}6$ s; \ $P=100$ W.
\end{svolto}

\section{Onde EM e ottica}
\begin{tab}{P{2.6cm}|L}
\hd Fenomeno & Formula\\\hline
Onda EM & $c=\frac1{\sqrt{\varepsilon_0\mu_0}}$, \ $v=\frac{c}{\sqrt{\varepsilon_r\mu_r}}$, \ $E=cB$, \ $I=\frac12c\varepsilon_0E_0^2$\\
Fotoni & $E=hf$; \ raggi X $\lambda_{min}=\frac{hc}{eV}$; \ Wien $\lambda_{max}T=0{,}29$ cm\,K\\
Snell & $n_1\sin\theta_1=n_2\sin\theta_2$; \ limite $\sin\theta_L=\frac{n_2}{n_1}$; \ $n=\sqrt{\varepsilon_r}$\\
Lastra & $s=\frac{d\sin(i-r)}{\cos r}$, \ cammino ottico $L=\frac{nd}{\cos r}$\\
Prisma (min. dev.) & $r=\frac\alpha2$, \ $n=\frac{\sin\frac{\alpha+\delta}{2}}{\sin\frac\alpha2}$\\
Young & max $d\sin\theta=m\lambda$, \ $\Delta y=\frac{\lambda L}{d}$; \ lamina $(n-1)s=m\lambda$\\
Reticolo & $b\sin\theta=m\lambda$, \ $R=mN$, \ $D=\frac{m}{b\cos\theta}$\\
Polarizzazione & Malus $I=I_0\cos^2\theta$; \ $\lambda/4$: $s(n_s-n_o)=\frac\lambda4$\\
\end{tab}

\section{Domande teoriche: immagine e punti da scrivere}
\begin{tab}{P{1.6cm}|P{2.5cm}|L}
\hd Domanda & Immagine & Punti\\\hline
Rotolamento (D1a) & compasso: il contatto è un perno fermo & $v=\omega R$; velocità $0,v,2v$; attrito statico senza lavoro; piano inclinato $a=\frac{g\sin\alpha}{1+I/MR^2}$\\
II principio (D1b) & mulino sulla cascata: serve la valle sotto & Kelvin, Clausius, equivalenza; Carnot; Clausius $\oint\delta Q/T\le0$; entropia\\
Gauss (D2a) & sacchetto chiuso: conta solo la carica dentro & flusso; enunciato; dimostrazione con la sfera; sfera, filo, piano; conduttori\\
Due fenditure (D2b) & due sassi nello stagno & Huygens; coerenza; $\delta=d\sin\theta$; max/min; $\Delta y$; $I=4I_0\cos^2$; inviluppo\\
Induzione & circuito conservatore (Lenz) & Faraday; $Blv$; esempi; autoinduzione\\
Maxwell & $E$ e $B$ si inseguono & 4 equazioni; corrente di spostamento; onde, $c$\\
I principio & conto in banca & $\Delta U=Q-W$; funzione di stato; trasformazioni; Mayer\\
Urti & passaggio della palla & $\vec p$ cost.; elastico/anelastico; CM\\
\end{tab}
\end{multicols*}
\end{document}
